%! Setting Up Determinants --- the 2 by 2 Determinant, Area, and Orientation — Unit 5, Lesson 5.0 — Lesson Plan
\documentclass[10pt]{article}
\usepackage{linalg-article}
\usepackage{linalg-boxes}
\usepackage{graphicx}   % -article does NOT load graphicx; needed for thumbnails/screenshots
\graphicspath{{images/}}
\newcommand{\TallMath}[1]{$\displaystyle #1\rule[-1.4em]{0pt}{3.2em}$}
\newcommand{\vv}{\mathbf{v}}
\newcommand{\ww}{\mathbf{w}}
\newcommand{\avec}{\mathbf{a}}
\newcommand{\bb}{\mathbf{b}}
\newcommand{\UnitNumberName}{Unit 5: Determinants and Linear Transformations \quad}
\newcommand{\LessonNumberName}{Lesson 5.0: Setting Up Determinants --- the 2 by 2 Determinant, Area, and Orientation}

\begin{document}
\begin{center}
    {\Large\bfseries \CourseName: \SchoolYear \\
    {\small\bfseries \UnitNumberName \LessonNumberName}}
\end{center}

\begin{tcolorbox}[colback=blush, colframe=burgundy, boxrule=0.9pt, arc=2mm,
    left=2mm, right=2mm, top=1.5mm, bottom=1.5mm]
    \textbf{Primary Objective:} Students can compute the \emph{determinant}
    $\det\begin{bsmallmatrix} a & b\\ c & d \end{bsmallmatrix}=ad-bc$ of a $2\times2$ matrix and
    read it \emph{geometrically}: its \emph{absolute value} is the \emph{area} of the parallelogram
    the columns span, and its \emph{sign} reports \emph{orientation} (a positive value keeps the
    plane's turning direction, a negative value \emph{flips} it). They can explain why
    $\det=0$ means the columns are \emph{parallel} --- the parallelogram collapses to a line, and the
    matrix is \emph{not invertible} (recall the Unit~2 inverse divides by $ad-bc$). They can describe
    where Unit~5 is going: $3\times3$ determinants and volume (5.1), properties and applications
    (5.2), and linear transformations, where $|\det|$ is the \emph{area-scaling factor} (5.3).
\end{tcolorbox}

\renewcommand{\arraystretch}{1.4}
\begin{skillbox}[Priority Ideas \& Skills]{goldbox}
\begin{tabularx}{\linewidth}{@{}X|X@{}}
\textbf{Priority Ideas \& Skills} & \textbf{Key Understandings (the ``why'')} \\[2pt]
\hline
\vspace{-6pt}
\begin{itemize}[leftmargin=1.1em, nosep, after=\vspace{2pt}]
  \item Compute a $2\times2$ determinant $ad-bc$.
  \item Read $|ad-bc|$ as the \textbf{area} of the parallelogram spanned by the columns.
  \item Read the \textbf{sign} as orientation: $+$ keeps, $-$ flips (a reflection); swapping columns negates it.
  \item Use $\det=0$: columns parallel $\Leftrightarrow$ zero area $\Leftrightarrow$ \emph{not invertible}.
\end{itemize}
&
\vspace{-6pt}
\begin{itemize}[leftmargin=1.1em, nosep, after=\vspace{2pt}]
  \item Four numbers collapse to \emph{one} that measures how the columns fill the plane.
  \item Determinant $=$ signed area: it turns a \emph{picture} (a parallelogram) into a \emph{number}.
  \item A negative sign is a \emph{mirror flip} --- the plane's turning direction reverses.
  \item $\det=0$ is the fingerprint of dependence: the box goes flat, the inverse dies.
\end{itemize}
\end{tabularx}
\end{skillbox}

\begin{skillbox}[{Vocabulary, Concepts, \& Theorems}]{greenbox}
\renewcommand{\arraystretch}{1.3}
\begin{tabularx}{\linewidth}{@{}l X@{}}
\textbf{Determinant $\det A$ (2$\times$2)} & For $A=\begin{bsmallmatrix} a & b\\ c & d \end{bsmallmatrix}$, the number $\det A=ad-bc$. \\
\textbf{Signed area} & $|\det A|$ is the area of the parallelogram spanned by $A$'s two columns; the \emph{sign} adds orientation. \\
\textbf{Orientation} & Positive $=$ same turning direction as the axes (counterclockwise); negative $=$ \emph{flipped} (a reflection). \\
\textbf{Singular / invertible} (recall §2) & $A$ is invertible exactly when $\det A\ne0$; the inverse $\frac{1}{ad-bc}\begin{bsmallmatrix} d & -b\\ -c & a \end{bsmallmatrix}$ divides by the determinant. \\
\textbf{Area-scaling factor} (preview 5.3) & Applying $A$ to a region multiplies every area by $|\det A|$. \\
\end{tabularx}
\end{skillbox}

\begin{fixedskillbox}[Activate Prior Knowledge \& Spiral Review]{blush}
    The Warm-Up rehearses the three ideas this lesson builds on:
    \textbf{(1)} reading the two \emph{columns} of a $2\times2$ matrix as \emph{vectors} (Unit~1),
    \textbf{(2)} the \emph{area} of a rectangle (base $\times$ height), and
    \textbf{(3)} the \emph{denominator} $ad-bc$ of the $2\times2$ inverse from Unit~2 --- the very
    number that decides whether the inverse exists. Debrief item~3 aloud: that denominator
    \emph{is} today's determinant, and ``the inverse blows up when $ad-bc=0$'' is the fact the
    whole lesson turns on.
\end{fixedskillbox}

\begin{skillbox}[Hook]{blush}
    Put a company logo inside the unit square and \emph{stretch and shear} it with a matrix --- the
    square becomes a parallelogram. Ask: \textbf{``How much bigger did the logo get, and did it get
    mirror-flipped?''} Both answers are hiding in \emph{one} number built from the matrix's four
    entries: $ad-bc$. Its \emph{size} is the new area (the stretch factor); its \emph{sign} says
    whether the image was flipped. Name the payoff: a single number reads off both the area and the
    orientation of what a matrix does. Pose the unit question: \textbf{``How can four numbers collapse
    into one that measures area and detects a flip?''} That number --- the determinant --- runs the
    whole unit: volume in 5.1, its rules and uses in 5.2, and transformations in 5.3.
\end{skillbox}

\begin{skillbox}[Lesson]{blush}
\begin{multicols}{2}
\textbf{1. The 2$\times$2 determinant.} For $A=\begin{bsmallmatrix} a & b\\ c & d \end{bsmallmatrix}$,
define $\det A=ad-bc$: down-diagonal product minus up-diagonal product. Students have already met
it --- it is the denominator of the Unit~2 inverse. Worked:
$\det\begin{bsmallmatrix} 2 & 1\\ 1 & 3 \end{bsmallmatrix}=6-1=5$.

\textbf{2. Determinant $=$ area.} The two columns span a parallelogram; $|\det A|$ is its
\textbf{area}. Axis-aligned first: columns $\begin{bsmallmatrix} 3\\0 \end{bsmallmatrix},
\begin{bsmallmatrix} 0\\2 \end{bsmallmatrix}$ make a $3\times2$ box, area $6=\det$. Shear it to
$\begin{bsmallmatrix} 3\\0 \end{bsmallmatrix},\begin{bsmallmatrix} 1\\2 \end{bsmallmatrix}$: base~$3$,
height~$2$, still area $6=3\cdot2-0\cdot1$. A picture becomes a number.

\columnbreak

\textbf{3. Orientation: the sign.} A \emph{positive} determinant keeps the plane's turning direction
(second column counterclockwise from the first, like the axes); a \emph{negative} one \textbf{flips}
it --- a mirror reflection. \emph{Swapping the two columns negates the determinant}:
$\det\begin{bsmallmatrix} 1 & 3\\ 2 & 0 \end{bsmallmatrix}=-6$ has the same area as
$\begin{bsmallmatrix} 3 & 1\\ 0 & 2 \end{bsmallmatrix}$ but the opposite orientation.

\textbf{4. When $\det=0$: collapse \& invertibility.} If the columns are \emph{parallel}, the
parallelogram is flat --- zero area --- and $\det=0$. E.g.\ $\begin{bsmallmatrix} 2 & 4\\ 1 & 2
\end{bsmallmatrix}$: $4-4=0$. That is exactly when $A$ is \textbf{not invertible} (the Unit~2
inverse divides by $ad-bc$). \textbf{Road ahead:} \textbf{5.1} $3\times3$ determinants and volume;
\textbf{5.2} properties and applications; \textbf{5.3} linear transformations, where $|\det A|$ is
the factor every area is multiplied by.
\end{multicols}
\end{skillbox}

\begin{skillbox}[Active Monitoring]{redbox}
Circulate during the notes practice and Tier work. Watch for and correct:
\begin{itemize}[leftmargin=1.2em, nosep]
  \item sign slips in $ad-bc$ --- computing $ad+bc$, or subtracting in the wrong order ($bc-ad$);
  \item forgetting \emph{absolute value} for area (a negative determinant still means positive area);
  \item reading $\det=0$ as ``no answer'' rather than ``columns parallel / flat box / no inverse.''
\end{itemize}
Cold-call prompts: ``Which diagonal is positive?'' ``A determinant of $-6$ --- what is the area, and
what does the minus sign mean?'' ``If $\det=0$, what do the two columns look like?''
\end{skillbox}

\begin{skillbox}[Group Work \& Differentiation]{redbox}
\textbf{Area, Orientation, and the Vanishing Box} activity (tiered; mirrors the notes).
\begin{multicols}{3}
\textbf{Tier R --- Remediate.} Compute $2\times2$ determinants; find the area of a parallelogram from
its columns; read whether a box is flat ($\det=0$) or not.

\columnbreak
\textbf{Tier A --- Approaching.} Use the \emph{sign} for orientation and confirm that swapping
columns flips it; solve for a missing entry that makes $\det=0$, and connect it to
\emph{invertibility} (recall §2).

\columnbreak
\textbf{Tier E --- Extension.} \emph{First look at transformations:} apply a matrix to the unit
square, find the area of the image, and see that it equals $|\det A|$ --- the area-scaling factor of
Lesson~5.3 (including a reflection with $\det<0$).
\end{multicols}
\end{skillbox}

\begin{skillbox}[Individual Work \& Assessment]{redbox}
\textbf{Exit Ticket} (independent, no notes): compute a $2\times2$ determinant; find the area of a
parallelogram and state whether its orientation is kept or flipped; and a
\textbf{conceptual/justification check} --- solve for the entry that makes $\det=0$ and explain what
that says about the columns and about invertibility. Collect and sort into ``got it / sign slip /
concept slip'' to decide whether 5.1 opens with a quick determinant re-warm.
\end{skillbox}

\begin{skillbox}[Reinforcement \& Extension]{goldbox}
\textbf{Homework:} compute determinants; areas of parallelograms from columns; orientation and the
swap rule; solve-for-the-entry to force $\det=0$ and tie it to invertibility (§2); and a short
justification item. \textbf{Extension:} apply a matrix to the unit square and show the image area is
$|\det A|$ --- a first taste of Lesson~5.3, including a reflection. \textbf{Preview:} Lesson~5.1,
\emph{$3\times3$ Determinants} --- the same idea one dimension up, where $|\det|$ becomes the
\emph{volume} of a box.
\end{skillbox}
\end{document}
